Agentic LLM applications often explore multiple candidate continuations from a shared context, while retaining at most one for subsequent steps.
In prefill--decode (PD) disaggregated serving, candidates submitted as follow-up requests may incur residual prefill and KV transfer before being discarded.
Existing decode-time pruning reduces subsequent generation, while KV-cache optimization avoids redundant computation, but neither fully addresses the early-stage costs of low-value candidates.
We propose \sys, a cross-stage pruning system that coordinates candidate admission, decode-side evaluation, and KV-cache management across the PD execution pipeline.
Specifically, (i) a lightweight admission policy uses residual scores and KV residency on both instances to reject low-value candidates before residual prefill and reuse cached KV; (ii) a decode-side probe further evaluates uncertain candidates through short continuations, reducing premature rejection while preserving pruning opportunities; and (iii) a copy-on-write context tree manages KV sharing and reclamation as branches are pruned or committed, retaining the selected path across agent steps.
Experiments on physically disaggregated vLLM instances show that admission reduces computed residual prefill by $23\%$ and ESC execution time by $12\%$, decode-side probing restores candidate coverage from $13/32$ to $20/32$ at $0.70\times$ the full-fan-out runtime, and copy-on-write management reduces median peak KV occupancy to $0.72\times$ that of automatic prefix caching.
